% Do not change document class, margins, fonts, etc.
\documentclass[a4paper,oneside,bibliography=totoc]{scrbook}


\usepackage[utf8]{inputenc}
\usepackage{graphicx}
\usepackage{latexsym}
\usepackage{amsmath}
\usepackage{amssymb}
\usepackage{tabularx}
\usepackage{booktabs}
\usepackage{algorithm}
\counterwithin{algorithm}{chapter}
\usepackage{algorithmic}
\usepackage{csquotes}
\renewcommand{\algorithmiccomment}[1]{\hfill\textit{// #1}}
\usepackage[usenames,dvipsnames,table]{xcolor}
\usepackage{array}
\usepackage{longtable}
\usepackage{xurl}
\usepackage[colorlinks,citecolor=Green]{hyperref}

\newcolumntype{L}[1]{>{\raggedright\arraybackslash\hspace{0pt}}p{#1}}
\newcolumntype{C}[1]{>{\centering\arraybackslash}p{#1}}
\definecolor{TableHeader}{RGB}{0,51,102}
\definecolor{LightTable}{RGB}{235,241,248}
\definecolor{Missing}{RGB}{192,0,0}
\newcommand{\tablegroup}[1]{\textbf{#1}\ }
\setlength{\emergencystretch}{3em}
\newcommand{\miss}[1]{#1\textsuperscript{\textcolor{Missing}{**}}}
\usepackage{lipsum}


\usepackage{natbib}
\bibliographystyle{chicagoa}
\setcitestyle{authoryear,round,semicolon,aysep={},yysep={,}} \let\cite\citep


\newtheorem{definition}{Definition} \newtheorem{proposition}{Proposition}

\begin{document}


\makeatletter
\renewcommand\thesection{\arabic{section}}
\renewcommand\thesubsection{\arabic{section}.\arabic{subsection}}
\renewcommand\thesubsubsection{\arabic{section}.\arabic{subsection}.\arabic{subsubsection}}
\makeatother

\subject{Web Data Integration \\ Final Report}
\title{Pharmaceutical Drug Integration: Linking US, EU and Canadian Approval Data}
\author{
  Gabriel Himmelein (1649181)\\
  Özlem Ölçer (2265825)\\
  Anatole Lobenko (2187322)\\
  Hai Nguyen (2267317)\\
  Thu Huyen Trang Trinh (2111810)
} \date{October 29, 2026}
\publishers{{\small Submitted to}\\
  Data and Web Science Group\\
  Dr. Ralph Peeters\\
  Prof. Dr. Christian Bizer\\
  University of Mannheim\\}
\maketitle

% ============================================================
\section{Use Case}
% ============================================================

This project builds a unified view of approved pharmaceutical drugs across the United States, European Union, and Canada by integrating three official regulatory datasets: Drugs@FDA (US), EMA Medicines (EU), and the Health Canada Drug Product Database (DPD). The central entity is a medicine identified by its brand name and active substance(s); each integrated record includes licence holder, regulatory status, product class, ATC code, dosage form, route, strengths, substance count, regional dates, and generic/biosimilar/orphan flags.

% ============================================================
\section{Requirements Suitability}
% ============================================================

\subsection{Schema and basic profile}

\begin{table}[htbp]
\centering
\caption{Overview of the input datasets}
\label{tab:datasets}
\scriptsize
\setlength{\tabcolsep}{4pt}
\begin{tabular}{|l|l|l|r|r|r|}
\hline
\textbf{Dataset} & \textbf{Source} & \textbf{Format} & \textbf{Rows} & \textbf{Entities} & \textbf{Attr.} \\
\hline
Drugs@FDA  & U.S. FDA    & TXT/TSV  & 51,839 &  8,175 & 56 \\
EMA Med.   & EMA         & XLSX/CSV &  2,746 &  2,620 & 39 \\
DPD        & Health Can. & TXT/CSV  & 42,042 & 41,636 & 87 \\
\hline
\end{tabular}
\end{table}

Table~\ref{tab:datasets} summarises the three input datasets. Drugs@FDA contains 51,839 product rows forming 8,175 entities across 29,377 applications (78.3\% ANDAs). EMA Medicines contains 2,746 rows forming 2,620 entities, predominantly human medicines (85.6\%). The Health Canada DPD contains 58,137 product records consolidated into 41,636 entities. High missingness in all three datasets is largely structural: optional fields, category-specific attributes, or regulatory-event dates that apply only to specific statuses.

\subsection{Integrated schema and overlap with input schemata}

The integrated schema contains 16 attributes (Table~\ref{tab:schema}); eleven occur in at
least two sources. \texttt{activeSubstances} and \texttt{strengths} are list-valued, while
the remaining 14 attributes are single-valued (including the numeric attribute
\texttt{activeSubstanceCount}). Strengths are represented as numeric value--unit pairs.

\begin{table}[htbp]
\centering
\caption{Attribute Intersection with Integrated Schema}
\label{tab:schema}
\scriptsize
\begin{tabular}{|l|l|>{\raggedright\arraybackslash}p{5.5cm}|}
\hline
\textbf{Attribute name} & \textbf{Attribute type} & \textbf{Datasets in which the attribute is found} \\
\hline
\texttt{brandName} & string & Drugs@FDA, EMA Medicines, DPD \\
\texttt{activeSubstances} & list (string) & Drugs@FDA, EMA Medicines, DPD \\
\texttt{activeSubstanceCount} & numeric (integer) & Drugs@FDA (derived), EMA Medicines (derived), DPD \\
\texttt{licenceHolder} & string & Drugs@FDA, EMA Medicines, DPD \\
\texttt{regulatoryStatus} & string (categorical) & Drugs@FDA, EMA Medicines, DPD \\
\texttt{productClass} & string (categorical) & Drugs@FDA (derived), EMA Medicines, DPD \\
\texttt{atcCode} & string & EMA Medicines, DPD \\
\texttt{generic} & boolean & Drugs@FDA (derived), EMA Medicines \\
\texttt{biosimilar} & boolean & EMA Medicines \\
\texttt{orphanMedicine} & boolean & EMA Medicines \\
\texttt{dosageForm} & string & Drugs@FDA, DPD \\
\texttt{route} & string & Drugs@FDA, DPD \\
\texttt{strengths} & list (numeric value, unit) & Drugs@FDA, DPD \\
\texttt{usApprovalDate} & date & Drugs@FDA \\
\texttt{euAuthorisationDate} & date & EMA Medicines \\
\texttt{caStatusDate} & date & DPD \\
\hline
\end{tabular}
\end{table}

% ============================================================
\subsection{Entity Identification and Count Requirement}
% ============================================================

A medicine is identified by its \texttt{brandName} and \texttt{activeSubstances}. Contrary
to initial assumptions, \texttt{licenceHolder} cannot be used as an identifying attribute
because the same medicine is frequently marketed by different regional subsidiaries or
unrelated licensing partners across jurisdictions (for example, \textit{Verkazia} is held
by \texttt{HARROW EYE} in the US, \texttt{Santen Oy} in the EU, and \texttt{APOTEX INC} in
Canada).

To estimate how many medicines appear in multiple datasets based on exact matching, we build
a normalized key (\texttt{medicine\_id}) from \texttt{brandName} and \texttt{activeSubstances}
by converting text to upper case, removing accents and dosage strengths, standardizing
punctuation, and sorting multiple active substances. Let $F$, $E$ and $D$ denote the sets of
\texttt{medicine\_id} values in Drugs@FDA, EMA Medicines and the DPD. The pairwise exact
overlaps are $|F \cap E| = 428$, $|F \cap D| = 949$ and $|E \cap D| = 594$, yielding a
deduplicated union of $\left|(F \cap E) \cup (F \cap D) \cup (E \cap D)\right| = 1{,}407$
distinct medicines -- well above the required minimum of 1,000.

However, exact matching misses many identical medicines because both identifying attributes
can vary across countries (Table~\ref{tab:difficult_match}):
\begin{itemize}
  \item \textbf{Active substance naming and salt forms:} The same substance can have different
  regional spellings (e.g., \texttt{CYCLOSPORINE} in the US and Canada vs.\
  \texttt{ciclosporin} in the EU for \textit{Verkazia}) or different chemical and salt
  representations (e.g., \texttt{GLYCOPYRROLATE} in the US, \texttt{glycopyrronium bromide}
  in the EU, and \texttt{GLYCOPYRRONIUM (GLYCOPYRRONIUM BROMIDE)} in Canada for
  \textit{Seebri}).
  \item \textbf{Brand name suffixes:} The same medicine may carry different device or
  formulation suffixes across markets, such as \texttt{SEEBRI NEOHALER} in Drugs@FDA versus
  \texttt{Seebri Breezhaler} in EMA Medicines and the DPD.
\end{itemize}

\begin{table}[htbp]
\centering
\caption{Examples of identical medicines with differing identifying attributes across sources}
\label{tab:difficult_match}
\footnotesize
\begin{tabularx}{\textwidth}{|l|l|X|}
\hline
\textbf{Dataset} & \textbf{Brand Name} & \textbf{Active Substance} \\
\hline
Drugs@FDA & \texttt{VERKAZIA} & \texttt{CYCLOSPORINE} \\
EMA Medicines & \texttt{Verkazia} & \texttt{ciclosporin} \\
Health Canada DPD & \texttt{VERKAZIA} & \texttt{CYCLOSPORINE} \\
\hline
Drugs@FDA & \texttt{SEEBRI NEOHALER} & \texttt{GLYCOPYRROLATE} \\
EMA Medicines & \texttt{Seebri Breezhaler} & \texttt{glycopyrronium bromide} \\
Health Canada DPD & \texttt{SEEBRI BREEZHALER} & \texttt{GLYCOPYRRONIUM (GLYCOPYRRONIUM BROMIDE)} \\
\hline
\end{tabularx}
\end{table}

% ============================================================
\section{Phase I: Data Selection and Translation}
% ============================================================

Phase~I translates three pharmaceutical datasets into a shared XML target
schema using Altova MapForce 2027. The outputs are validated XML files
consumed by Phases~II and~III.

% ------------------------------------------------------------
\subsection{Target Schema}
% ------------------------------------------------------------

The target schema (\texttt{medicines.xsd}) represents one medicine entity
per source record, identified by \texttt{brandName} and
\texttt{activeSubstances}, with 16 attributes and two metadata fields
(\texttt{id}, \texttt{source}). Three conventions apply across all
mappers: missing values are omitted rather than written as empty tags;
repeating values use typed wrapper elements (\texttt{activeSubstances},
\texttt{strengths}); and all dates follow ISO~8601 (YYYY-MM-DD).
Controlled vocabularies are enforced for \texttt{regulatoryStatus} (six
values: \textit{Authorised, Suspended, Pending, Withdrawn, Application
withdrawn, Refused}) and \texttt{productClass} (\textit{Human,
Veterinary, Disinfectant}).

% ------------------------------------------------------------
\subsection{Shared Helper Functions}
% ------------------------------------------------------------

A shared MapForce library (\texttt{wdi\_helpers.mfd}) provides ten
user-defined functions imported by all mappers
(Table~\ref{tab:udfs}). Lookup XML files cover 25 regulatory statuses,
227 dosage forms, 147 routes, and 202 strength units across all three
sources.

\begin{table}[htbp]
\centering
\caption{Shared UDFs used across mappers}
\label{tab:udfs}
\scriptsize
\begin{tabular}{|l|>{\raggedright\arraybackslash}p{7.5cm}|}
\hline
\textbf{Function} & \textbf{Purpose} \\
\hline
\texttt{makeMedicineId} & Constructs a stable entity ID from source prefix and record key \\
\texttt{getActiveSubstances} & Tokenizes substance string; supports INN fallback for blank primary field \\
\texttt{standardizeRegulatoryStatus} & Maps raw status values to the six controlled vocabulary entries \\
\texttt{standardizeProductClass} & Maps category values to Human / Veterinary / Disinfectant \\
\texttt{standardizeAtcCodes} & Combines human and vet ATC codes; deduplicates, sorts, corrects typos \\
\texttt{yesNoToBoolean} & Converts Yes/No strings to \texttt{xs:boolean}; omits element when false \\
\texttt{parseAuthorisationDate} & Converts DD/MM/YYYY to ISO~8601 YYYY-MM-DD \\
\texttt{stdDosageForm} / \texttt{dosageFormList} & Standardizes single or multi-value dosage form strings \\
\texttt{stdRoute} / \texttt{routeList} & Standardizes single or multi-value administration routes \\
\texttt{splitStrengths} & Tokenizes a DPD strength string into (value, unit) pairs \\
\hline
\end{tabular}
\end{table}

% ------------------------------------------------------------
\subsection{EMA Medicines Mapper}
% ------------------------------------------------------------

The EMA dataset is a single flat CSV table of 2,746 rows, with no
pre-joining or consolidation required. Each row maps directly to one
output \texttt{<medicine>} element. Because EMA does not publish strength
or dosage form data, \texttt{dosageForm}, \texttt{route}, and
\texttt{strengths} are structurally absent from EMA output.
Table~\ref{tab:ema_mapping} lists all field correspondences.

\begin{table}[htbp]
\centering
\caption{EMA Medicines field mapping}
\label{tab:ema_mapping}
\scriptsize
\begin{tabular}{|l|l|l|}
\hline
\textbf{Target field} & \textbf{Source column} & \textbf{Transformation} \\
\hline
\texttt{id}                   & EMA product number           & \texttt{makeMedicineId("EMA", number)} \\
\texttt{brandName}            & Name of medicine             & Direct copy \\
\texttt{activeSubstances}     & Active substance / INN       & \texttt{getActiveSubstances}; INN fallback \\
\texttt{activeSubstanceCount} & (derived)                    & \texttt{count(activeSubstances)} \\
\texttt{licenceHolder}        & MA holder                    & Direct copy \\
\texttt{regulatoryStatus}     & Medicine status              & \texttt{standardizeRegulatoryStatus} \\
\texttt{productClass}         & Category                     & \texttt{standardizeProductClass} \\
\texttt{atcCode}              & ATC code (human) + ATCvet    & \texttt{standardizeAtcCodes} \\
\texttt{generic} / \texttt{biosimilar} / \texttt{orphanMedicine} & Generic / Biosimilar / Orphan & \texttt{yesNoToBoolean} \\
\texttt{euAuthorisationDate}  & Marketing authorisation date & \texttt{parseAuthorisationDate} \\
\hline
\end{tabular}
\end{table}

The most intricate correspondence handles active substances. For 57 of
the 2,746 rows, both \texttt{Active substance} and the \texttt{INN /
common name} column are blank; because \texttt{activeSubstances} is
mandatory in the target schema, these rows are excluded via a row filter.
A further 28 rows have a blank primary field but a populated INN column:
the mapper applies an INN fallback, normalizing whitespace and selecting
the non-empty value. Commas and the word ``and'' are not treated as
separators, preserving multi-word biological descriptions. The result is
\textbf{2,689 output records}.

EMA maintains separate human and veterinary ATC code columns.
\texttt{standardizeAtcCodes} concatenates both, strips asterisks used
for provisional codes, drops the \textit{Not yet assigned} placeholder,
corrects a known EMA typo (\texttt{VO4D} $\rightarrow$ \texttt{V04D}),
then deduplicates, sorts, and joins codes with \texttt{"; "}. Date
conversion (\texttt{parseAuthorisationDate}) reorders DD/MM/YYYY to
YYYY-MM-DD. The three boolean flags are stored as Yes/No strings;
\texttt{yesNoToBoolean} converts them and omits the element when false,
consistent with the \texttt{minOccurs="0"} schema convention. The full
export passed XSD validation with 2,689 unique \texttt{EMA\_} IDs and no
errors.

% ------------------------------------------------------------
\subsection{Drugs@FDA Mapper}
% ------------------------------------------------------------

The Drugs@FDA mapper specification is defined in the project repository;
implementation is pending.

% ------------------------------------------------------------
\subsection{Health Canada DPD Mapper}
% ------------------------------------------------------------

The DPD source is a pre-consolidated CSV of \textbf{41,636 medicine
entities} derived from 58,137 product records across 12 Health Canada
tables (\texttt{dpd\_entities.csv}). Unlike the EMA flat table, each DPD
entity already aggregates multiple product records, so the mapper
operates on one row per entity without further consolidation. The DPD
source provides all 16 target attributes including \texttt{dosageForm},
\texttt{route}, and \texttt{strengths}, making it the most complete of
the three sources.

\begin{table}[htbp]
\centering
\caption{Health Canada DPD field mapping}
\label{tab:dpd_mapping}
\scriptsize
\begin{tabular}{|l|l|l|}
\hline
\textbf{Target field} & \textbf{Source column} & \textbf{Transformation} \\
\hline
\texttt{id}                   & medicine\_id          & \texttt{makeMedicineId("DPD", id)} \\
\texttt{brandName}            & brand\_name           & Direct copy \\
\texttt{activeSubstances}     & active\_substances    & \texttt{getActiveSubstances}; separator \texttt{"; "} \\
\texttt{activeSubstanceCount} & (derived)             & \texttt{count(activeSubstances)} \\
\texttt{licenceHolder}        & licence\_holder       & Direct copy \\
\texttt{regulatoryStatus}     & current\_status       & \texttt{standardizeRegulatoryStatus} \\
\texttt{productClass}         & product\_class        & \texttt{standardizeProductClass} \\
\texttt{atcCode}              & atc\_code             & \texttt{standardizeAtcCodes} \\
\texttt{biosimilar}           & biosimilar            & \texttt{yesNoToBoolean} \\
\texttt{dosageForm}           & dosage\_forms         & \texttt{dosageFormList} \\
\texttt{route}                & routes                & \texttt{routeList} \\
\texttt{strengths}            & strengths             & \texttt{splitStrengths} $\rightarrow$ (value, unit) pairs \\
\texttt{caStatusDate}         & status\_date          & Direct copy (ISO date) \\
\hline
\end{tabular}
\end{table}

The main mapping challenge in DPD is strength parsing. Each entity row
carries a semicolon-delimited strength string (e.g.\
\texttt{"10MG; 0.5\%"}); \texttt{splitStrengths} tokenizes this into
individual tokens, from which \texttt{strengthValue} extracts the numeric
component and \texttt{strengthUnit} normalizes the unit via
\texttt{lookups/strength\_unit.xml}. Tokens without a parseable numeric
value are filtered out before the \texttt{<strength>} wrapper is created,
and the wrapper itself is omitted when no valid tokens remain. Regulatory
status presents a similar aggregation challenge: DPD entities can carry
multiple historical statuses (e.g.\ \texttt{"MARKETED; CANCELLED POST
MARKET"}), so \texttt{standardizeRegulatoryStatus} selects the highest
active status according to the precedence hierarchy
$\text{Authorised} > \text{Suspended} > \cdots > \text{Refused}$.

The full export passed XSD validation with \textbf{41,636 records},
100,094 strength pairs across 41,574 medicines, and no schema errors.




% ============================================================
\section{Phase II: Identity Resolution}
% ============================================================




% ============================================================
\section{Phase III: Data Fusion}
% ============================================================




% ============================================================
\section{Conclusion}
% ============================================================




% ============================================================
% SOURCES
% ============================================================

\section*{Sources}
\addcontentsline{toc}{section}{Sources}

\begin{itemize}
    \item U.S. Food and Drug Administration,
    \emph{Drugs@FDA Data Files}.\\
    \url{https://www.fda.gov/drugs/drug-approvals-and-databases/drugsfda-data-files}

    \item European Medicines Agency,
    \emph{Download medicine data}.\\
    \url{https://www.ema.europa.eu/en/medicines/download-medicine-data}

    \item Health Canada,
    \emph{Drug Product Database}.\\
    \url{https://www.canada.ca/en/health-canada/services/drugs-health-products/drug-products/drug-product-database.html}
\end{itemize}

\end{document}
